\section*{Chapter \ref{ch:in:the-options-market-making-houses} --- The Options Market-Making Houses}

\begin{solution}{exo:in:the-options-market-making-houses:1}
$15\times40\times2=1\,200$ series.
\end{solution}

\begin{solution}{exo:in:the-options-market-making-houses:2}
$1\,158/2\,773=41.8\%$.
\end{solution}

\begin{solution}{exo:in:the-options-market-making-houses:3}
IMC (the Amsterdam equity options exchange), Susquehanna (the Philadelphia Stock Exchange), Belvedere Trading (the
Chicago Board Options Exchange).
\end{solution}

\begin{solution}{exo:in:the-options-market-making-houses:4}
$1\,369/1\,158-1=18.2\%$ in 2024 and $1\,769/1\,369-1=29.2\%$ in 2025. In 2024 net trading income grew 26\%: costs
(pay, technology, tax) grew faster than income, so the margin fell from 41.8\% to 39.2\%. The releases do not break
costs down, so which cost grew is not known from them.
\end{solution}

\begin{solution}{exo:in:the-options-market-making-houses:5}
$4\,556/2\,400=1.90$ to $4\,556/2\,000=2.28$: \euro1.90 to \euro2.28 million per head.
\end{solution}

\begin{solution}{exo:in:the-options-market-making-houses:6}
Its book is a surface of volatility, skew and time-decay exposures that can be hedged only partly and slowly, by
trading other options; and its quoting obligations make it take positions it would not choose.
\end{solution}

\begin{solution}{exo:in:the-options-market-making-houses:7}
30.4\% in 2024 and 34.0\% in 2025; $4\,556/2\,500=\text{\euro}1.82$ million per head.
\end{solution}

\begin{solution}{exo:in:the-options-market-making-houses:8}
One house that chooses to publish is not a sample of the industry: it may be among the largest and most successful;
its headcount is a lower bound, so its figure is an upper bound; and 2024--2025 were years of high volumes by the
house's own account. Per-head income varies by firm, by year and by what is counted.
\end{solution}

\begin{solution}{pb:in:the-options-market-making-houses:1}
\begin{enumerate}
\item A \omterm{def:in:a-map-of-the-industry:ptf}{principal trading firm} whose core business is quoting listed options across strikes and expiries and hedging
them; Book~1's options market maker is the role on a venue, with its obligations.
\item Amsterdam equity options exchange (IMC), Philadelphia Stock Exchange (Susquehanna), Chicago Board Options Exchange
(Belvedere Trading).
\item 1980s 4, 1990s 1, 2000s 1.
\item A firm founded by people who leave an existing firm to do the same business; no firm states such an origin on
its own pages.
\item It runs a hybrid market combining \omterm{def:in:proprietary-market-makers-and-high-frequency-firms:floor}{open outcry} floor trading with electronic trading.
\item $20\times50\times2=2\,000$.
\item The number of instruments, the risk held, the obligations, the model (a surface, not a price).
\item Pricing, the Greeks, surfaces, inventory and hedging; practice through games, simulated markets and supervised
work.
\item Hundreds of thousands of quotes move with every change in the underlying; no person can update them.
\item The employment contract.
\item Net trading income \euro2\,773, 3\,494, 4\,556 million; net profit \euro1\,158, 1\,369, 1\,769 million.
\item 26.0\% and 30.4\%.
\item 41.8\%, 39.2\%, 38.8\%.
\item 30.4\% and 34.0\%.
\item \euro1.66 million in 2024; at most \euro2.28 million in 2025.
\item \euro4\,556 million, 38.8\% and 34.0\%; three of the six firms (50\%) tie their founding to a floor.
\item More than 2\,000 employees means dividing by at least 2\,000.
\item It chose to publish, it is one firm, and its years were busy; there is no sample.
\item Its UK subsidiary's accounts, or its US broker-dealer's public statement (chapter~2).
\item The houses train for pricing and risk judgement on a surface, which is a different job from racing for a
price; one house's numbers show the business can be very profitable, not what a given house earns.
\end{enumerate}
\end{solution}

\begin{solution}{iq:in:the-options-market-making-houses:1}
Its inventory is a set of options whose risk (vega, skew, gamma) cannot be flattened quickly without paying the spread
on many series; it hedges delta continuously and carries the rest, managing it over days.
\iqlookfor{the difference between hedging delta and flattening a book of Greeks.}
\end{solution}

\begin{solution}{iq:in:the-options-market-making-houses:2}
$12\times30\times2=720$ series, two sides each: 1\,440 quotes.
\iqlookfor{calls and puts, bids and offers.}
\end{solution}

\begin{solution}{iq:in:the-options-market-making-houses:3}
For the appointment's benefits: fee rebates, allocation or priority advantages and the flow that comes with being
present (Book~1, chapter~24). The cost is being present in series and at times the firm would avoid, adversely selected
and with capital tied up.
\iqlookfor{obligations bought with privileges.}
\end{solution}

\begin{solution}{iq:in:the-options-market-making-houses:4}
Roughly $1.29/1.12\approx1.15$: return on equity up by about 15\%. Check whether the year was unusually volatile,
whether the growth came from one business or event, and whether risk taken rose.
\iqlookfor{ratio arithmetic and the cycle before skill.}
\end{solution}

\begin{solution}{iq:in:the-options-market-making-houses:5}
How the firm's own models and surfaces behave on its own flow; how to manage a real book under limits; how the
desk's tools and alerts work; the market's microstructure on specific venues.
\iqlookfor{practice, tools and judgement, not theory.}
\end{solution}

\begin{solution}{iq:in:the-options-market-making-houses:6}
Long gamma earns from moves but decays; short vega loses if implied volatility rises. Check the delta hedge, the
overnight event calendar and the vega by expiry; reduce the vega if an event could lift implied volatility, and
decide how much gamma to keep given the expected overnight move against the decay.
\iqlookfor{Greeks as risks with prices, and an overnight plan.}
\end{solution}
